\section*{अध्याय \ref{ch:g11:becoming-male-female} --- नर या मादा बनना}

\begin{solution}{exo:g11:becoming-male-female:1}
पिता का: हर अंडाणु एक X ढोता है, जबकि आधे शुक्राणु X ढोते हैं और आधे Y,
इसलिए XX और XY बराबर संभावना वाले हैं।
\end{solution}

\begin{solution}{exo:g11:becoming-male-female:2}
दो अविभेदित जननग्रंथियाँ, और हर एक के बग़ल में दो वाहिनियाँ: एक वुल्फ़ियन
वाहिनी (भावी नर तंत्र) और एक म्यूलेरियन वाहिनी (भावी मादा तंत्र)। XX और XY
भ्रूणों में एक जैसी।
\end{solution}

\begin{solution}{exo:g11:becoming-male-female:3}
Y \omterm{def:g11:cell-cycle-mitosis:chromosome}{गुणसूत्र} पर; हफ़्ता 7 के आसपास अविभेदित जननग्रंथि में अभिव्यक्त होकर उसका
\omterm{def:g11:gene-expression:protein}{प्रोटीन} उन \omterm{def:g10:universal-dna:gene}{जीनों} को चालू कर देता है जो उस जननग्रंथि को वृषण बना देते हैं।
\end{solution}

\begin{solution}{exo:g11:becoming-male-female:4}
\omterm{prop:g11:becoming-male-female:hormones}{टेस्टोस्टेरोन} वुल्फ़ियन वाहिनियों को बनाए रखता है और उन्हें नर तंत्र में
विकसित करता है; और प्रति-म्यूलेरियन हॉर्मोन म्यूलेरियन वाहिनियों को सिकोड़
देता है।
\end{solution}

\begin{solution}{exo:g11:becoming-male-female:5}
मस्तिष्क का एक केंद्र किसी हॉर्मोन की धड़कनें छोड़ता है, जो पीयूष ग्रंथि से
LH और FSH स्रावित कराती हैं, और वे जननग्रंथियाँ सक्रिय कर देते हैं। गौण
लैंगिक लक्षण: लैंगिक हॉर्मोनों के पैदा किए बदलाव --- शरीर के बाल, आवाज़,
स्तन, पेशी तथा चर्बी का बँटवारा, वृद्धि-उछाल, और जनन अंगों की वृद्धि।
\end{solution}

\begin{solution}{exo:g11:becoming-male-female:6}
अपने Y \omterm{def:g11:cell-cycle-mitosis:chromosome}{गुणसूत्र} के बावजूद उसमें मादा तंत्र विकसित हुआ: यानी नर तंत्र को
वृषण के हॉर्मोन चाहिए, और मादा तंत्र बिना किसी जननग्रंथि के बन जाता है।
\end{solution}

\begin{solution}{exo:g11:becoming-male-female:7}
\omterm{prop:g11:becoming-male-female:hormones}{टेस्टोस्टेरोन} नर वाहिनियाँ बनाता तो है पर मादा वाली हटाता नहीं; उसके लिए
वृषण को कोई दूसरा पदार्थ स्रावित करना ही होगा --- यानी AMH।
\end{solution}

\begin{solution}{exo:g11:becoming-male-female:8}
वृषण (Sry काम करता है); वुल्फ़ियन वाहिनियाँ बचीं, म्यूलेरियन सिकुड़ीं; और
बाहरी रचना नर। बंध्य: शुक्राणु बनाने के लिए Y के वे \omterm{def:g10:universal-dna:gene}{जीन} चाहिए जो वह पारजीन
देता ही नहीं।
\end{solution}

\begin{solution}{exo:g11:becoming-male-female:9}
एस्ट्रोजन: लगभग 11 पर 50\%, और पठार लगभग 14--15 पर। \omterm{prop:g11:becoming-male-female:hormones}{टेस्टोस्टेरोन}: लगभग
13 पर 50\%, और पठार लगभग 16--17 पर।
\end{solution}

\begin{solution}{exo:g11:becoming-male-female:10}
हड्डियाँ तब बढ़ना बंद कर देती हैं जब लैंगिक हॉर्मोन अपने वयस्क स्तर पर
पहुँच जाए; और एस्ट्रोजन \omterm{prop:g11:becoming-male-female:hormones}{टेस्टोस्टेरोन} से लगभग दो साल पहले वहाँ पहुँच जाते
हैं, इसलिए लड़कियों की वृद्धि पट्टिकाएँ पहले बंद हो जाती हैं।
\end{solution}

\begin{solution}{exo:g11:becoming-male-female:11}
कोई वृषण नहीं: वह जननग्रंथि अंडाशय जैसी बन जाती है; न \omterm{prop:g11:becoming-male-female:hormones}{टेस्टोस्टेरोन}, न
AMH: म्यूलेरियन वाहिनियाँ विकसित, वुल्फ़ियन सिकुड़ीं, और जन्म पर बाहरी रचना
मादा। \omterm{prop:g11:becoming-male-female:puberty}{यौवनारंभ} पर वे जननग्रंथियाँ, जिनमें सामान्य अंडाणु नहीं हैं, कम ही
स्रावित करती हैं; और हॉर्मोन के इलाज के बिना गौण लक्षण ठीक से विकसित नहीं
होते।
\end{solution}

\begin{solution}{exo:g11:becoming-male-female:12}
वृषण (\omterm{prop:g11:becoming-male-female:sry}{SRY} काम करता है)। वुल्फ़ियन वाहिनियाँ सिकुड़ जाती हैं (\omterm{prop:g11:becoming-male-female:hormones}{टेस्टोस्टेरोन}
का कोई जवाब नहीं); म्यूलेरियन वाहिनियाँ भी सिकुड़ जाती हैं (AMH काम करता
है): यानी किसी भी तरह का भीतरी तंत्र नहीं। बाहरी रचना मादा (उसके नर रूप के
लिए \omterm{prop:g11:becoming-male-female:hormones}{टेस्टोस्टेरोन} की क्रिया चाहिए)। और \omterm{prop:g11:becoming-male-female:puberty}{यौवनारंभ} पर \omterm{prop:g11:becoming-male-female:hormones}{टेस्टोस्टेरोन} से बने
एस्ट्रोजन स्तन तथा मादा गठन देते हैं, पर शरीर के बाल नहीं।
\end{solution}

\begin{solution}{exo:g11:becoming-male-female:13}
कोई जननग्रंथि एक बार वृषण या अंडाशय बन जाए, तो उसके हॉर्मोन वाहिनियों और
शरीर को ठीक स्तनधारियों जैसा ही गढ़ सकते हैं; फ़र्क़ पहले क़दम का है ---
जननग्रंथि का फ़ैसला करने वाले स्विच के रूप में \omterm{prop:g11:becoming-male-female:sry}{SRY} की जगह कोई ताप-संवेदी
प्रक्रिया आ जाती है।
\end{solution}

\begin{solution}{exo:g11:becoming-male-female:14}
X सैकड़ों ऐसे \omterm{def:g10:universal-dna:gene}{जीन} ढोता है जो हर \omterm{def:g10:cells-common-unit:cell}{कोशिका} को चाहिए, इसलिए उसके बिना कोई भ्रूण
बचता ही नहीं; और Y वृषण तय करने वाले तथा शुक्राणु से जुड़े \omterm{def:g10:universal-dna:gene}{जीनों} के अलावा
लगभग कुछ नहीं ढोता, इसलिए XX व्यक्तियों को ज़रूरी कोई चीज़ नहीं खटकती। Y
केवल बेटों तक इसलिए जाता है कि बेटियों को पिता का X मिलता है।
\end{solution}

\begin{solution}{exo:g11:becoming-male-female:15}
उसका मतलब है कि जब वृषण का कोई हॉर्मोन दख़ल न दे तब मादा तंत्र विकसित हो
जाता है --- यानी यह इस बारे में एक तथ्य है कि कौन-से संकेत चाहिए, कोई
क्रम-निर्धारण नहीं। वह किसी साझा भ्रूणीय योजना को दर्शाता है, जिस पर
\omterm{def:g10:cells-common-unit:cell}{कोशिकाओं} का एक वंश, यानी वृषण, दो हॉर्मोनों से एक विकल्प थोप देता है:
दोनों लिंग एक ही कार्यक्रम के दो नतीजे हैं, जो किसी एक स्विच से अलग पड़ते
हैं।
\end{solution}

\begin{solution}{pb:g11:becoming-male-female:1}
\textbf{1.} मादा तंत्र के लिए अंडाशय की ज़रूरत नहीं, वह बिना किसी
जननग्रंथि के बन जाता है; और नर तंत्र के लिए तथा मादा वाहिनियाँ हटाने के लिए
वृषण चाहिए।

\textbf{2.} 22वें दिन के बाद वाहिनियाँ जननग्रंथि के असर में विभेदित होना
शुरू कर चुकी होती हैं; और उस प्रयोग को फ़ैसले से पहले काम करना ही था।

\textbf{3.} वृषण के हॉर्मोन, जिन्हें रक्त वाहिनियों तक ढोता है: वह
प्रत्यारोपण, वाहिनियों से दूर होते हुए भी, उन्हें नर बना देता है।

\textbf{4.} \omterm{prop:g11:becoming-male-female:hormones}{टेस्टोस्टेरोन} नर वाहिनियाँ बनाता है; वह मादा वाली हटाता नहीं।

\textbf{5.} वृषण का कोई दूसरा पदार्थ म्यूलेरियन वाहिनियाँ सिकोड़ता ही
होगा; और उसे वृषणों से निकालकर अकेले बिठाकर पक्का किया जा सकता था, इस
अपेक्षा के साथ कि मादा वाहिनियाँ ग़ायब हो जाएँ और नर वाली बनें नहीं ---
और AMH ठीक यही करता है।

\textbf{6.} जिस ओर शल्यक्रिया हुई वहाँ न \omterm{prop:g11:becoming-male-female:hormones}{टेस्टोस्टेरोन}, न AMH: यानी मादा
वाहिनी, नर वाहिनी नहीं। और सही-सलामत ओर: नर वाहिनी, और मादा वाहिनी
सिकुड़ी हुई।

\textbf{7.} प्रत्यारोपण वाली ओर नर वाहिनी और कोई मादा वाहिनी नहीं; और
दूसरी ओर केवल मादा वाहिनी।

\textbf{8.} क्रिस्टल वाली ओर दोनों वाहिनियाँ मौजूद; और दूसरी ओर केवल मादा
वाहिनी।

\textbf{9.} हर शिशु ख़ुद अपना नियंत्रण बन जाता है: उसके दोनों ओर का फ़र्क़
केवल वृषण या \omterm{prop:g11:becoming-male-female:hormones}{टेस्टोस्टेरोन} के स्थानीय होने या न होने से ही आ सकता है।

\textbf{10.} उसके अंडाशय हटा दिए गए हैं या दबा दिए गए हैं, और किसी दूसरे
जानवर से आया वह प्रत्यारोपित वृषण किसी पराए मेज़बान में शुक्राणु बनाता
नहीं; इसलिए तंत्र नर है पर कोई युग्मक बनते ही नहीं।

\textbf{11.} \omterm{prop:g11:becoming-male-female:sry}{SRY} मौजूद: वृषण; \omterm{prop:g11:becoming-male-female:hormones}{टेस्टोस्टेरोन} और AMH: नर तंत्र, नर रचना; और
\omterm{prop:g11:becoming-male-female:puberty}{यौवनारंभ} पर \omterm{prop:g11:becoming-male-female:hormones}{टेस्टोस्टेरोन}: नर लक्षण। बंध्य (शुक्राणु के लिए Y के \omterm{def:g10:universal-dna:gene}{जीन} नहीं)।

\textbf{12.} काम करता कोई \omterm{prop:g11:becoming-male-female:sry}{SRY} नहीं: कोई वृषण नहीं; कोई हॉर्मोन नहीं: मादा
तंत्र और मादा रचना; और जननग्रंथियाँ ठीक से विकसित नहीं।

\textbf{13.} वृषण; वुल्फ़ियन वाहिनियाँ विकसित (\omterm{prop:g11:becoming-male-female:hormones}{टेस्टोस्टेरोन} काम करता है);
म्यूलेरियन वाहिनियाँ भी टिकी रहती हैं (AMH का कोई जवाब नहीं): यानी ऐसा नर
जिसके पास इसके अलावा गर्भाशय और अंडवाहिनियाँ भी हैं; और नर रचना तथा नर
\omterm{prop:g11:becoming-male-female:puberty}{यौवनारंभ}।

\textbf{14.} केवल तीसरा: काम करते वृषण और नर वाहिनियाँ, हालाँकि टिका हुआ
गर्भाशय वृषणों के नीचे उतरने में अड़चन डाल सकता है। पहले के पास शुक्राणु
वाले \omterm{def:g10:universal-dna:gene}{जीन} नहीं; और दूसरे के पास काम करती जननग्रंथियाँ ही नहीं।

\textbf{15.} गुणसूत्रीय लिंग शारीरिक लिंग की भविष्यवाणी केवल बीच के क़दमों
से करता है: काम करता \omterm{prop:g11:becoming-male-female:sry}{SRY}, काम करते हॉर्मोन, काम करते ग्राही। किसी भी क़दम
को बदलिए और दोनों स्तर अलग हो जाते हैं।

\textbf{16.} लड़कियाँ: 10 पर 20\%, उछाल लगभग 11 से, और वृद्धि लगभग 15 पर
रुकती है। लड़के: लगभग 12.5 पर 20\%, उछाल लगभग 13.5 से, और लगभग 17 पर
रुकती है।

\textbf{17.} लड़कियों को लगभग 4 साल, लड़कों को लगभग 3.5 --- पर लड़के अधिक
ऊँचे आधार से शुरू करते हैं, क्योंकि वे दो साल अधिक बचपन की दर से बढ़ चुके
होते हैं।

\textbf{18.} लड़के लगभग $3.5 \times 8 = \qty{28}{cm}$, लड़कियाँ $4 \times
7 = \qty{28}{cm}$: यानी दोनों के उछाल मिलते-जुलते हैं; और वयस्क फ़र्क़
ज़्यादातर लड़कों के उछाल से पहले के बचपन की दो अतिरिक्त साल की वृद्धि से
आता है।

\textbf{19.} जल्दी आई चढ़ाई वृद्धि पट्टिकाओं को सालों पहले बंद करा देती
है: वह बच्चा अभी तेज़ बढ़ता है पर कम वयस्क लंबाई पर रुक जाता है। हॉर्मोन
टालने से बचपन की वृद्धि लंबी खिंच जाती है।

\textbf{20.} \omterm{prop:g11:becoming-male-female:hormones}{टेस्टोस्टेरोन}, जो नर वाहिनियाँ बनाता है, और
प्रति-म्यूलेरियन हॉर्मोन, जो मादा वाली हटा देता है; और दोनों वृषण बनाता
है, यानी वह अंग जिसे अकेला \omterm{def:g10:universal-dna:gene}{जीन} \omterm{prop:g11:becoming-male-female:sry}{SRY} चालू करता है।
\end{solution}
